\section{Introduction}\label{sec:intro}

Deep reinforcement learning (\drl) has had a massive impact on the
field of machine learning and has led to remarkable successes in the
solution of many challenging
tasks~\citep{mnih2015human,AlphaGo,silver2017mastering}.  While neural
networks have been shown to be very effective in learning good
policies, the expressivity of these models makes them difficult to
interpret or to be checked for consistency for some desired
properties, and casts a cloud over the use of such representations in
safety-critical applications.

Motivated to overcome this problem, we propose a learning framework,
called {\em Programmatically Interpretable Reinforcement Learning}
(\pirl)\footnote{\pirl is pronounced Pi-R-L (as in $\pi$-\rl)}, that is based on the idea of learning policies that are
represented in a human-readable language.  The \pirl framework is
parameterized on a high-level {\em programming language} for policies.
A problem instance in \pirl is similar to a one in traditional \rl,
but also includes a {\em (policy) sketch} that syntactically defines a set of programmatic policies
in this language.  The objective is to find a program in this set with
maximal long-term reward.

Intuitively, the policy programming language and the sketch characterize what we
consider ``interpretable''. In addition to interpretability, the
syntactic restriction on policies has three key benefits. First, the language
can be used to implicitly encode the learner's inductive bias that will
be used for generalization. Second, the language can allow effective
pruning of undesired policies to make the search for a good policy
more efficient. Finally, it 
allows us to use symbolic program verification techniques to formally reason
about the learned policies and check consistency with correctness
properties. At the same time, policies in \pirl can have rich 
semantics, for example allowing actions to depend on events far back
in history. 

A key technical challenge in \pirl is that the space of policies
permitted in an instance can be vast and nonsmooth, making
optimization extremely challenging. To address this, we propose a new
algorithm called {\em Neurally Directed Program Synthesis}
(\algo). The algorithm first uses \drl to compute a neural policy
network that has high performance, but may not be expressible in the
policy language. This network is then used to direct a local search
over programmatic policies. In each iteration of this search, we
maintain a set of ``interesting'' inputs, and update the program so as
to minimize the distance between its outputs and the outputs of the
neural policy (an ``oracle'') on these inputs. The set of interesting
inputs is updated as the search progresses. This strategy, inspired by
imitation learning~\cite{dagger,schaal1999imitation}, allows us to
perform direct policy search in a highly nonsmooth policy space.

We evaluate our approach in the task of learning to drive a simulated
car in the \torcs car-racing environment~\citep{TORCS}, as well as
three classic control games (we discuss the former in the main paper,
and the latter in the Appendix). Experiments demonstrate that \algo is
able to find interpretable policies that, while not as performant as
the policies computed by \drl, pass some significant performance
bars. Specifically, in \torcs, our policy sketch allows an unbounded
set of programs with branches guarded by unknown conditions, each
branch representing a Proportional-Integral-Derivative (PID)
controller~\citep{aastrom1995pid} with unknown parameters. The policy we
obtain can successfully
complete a lap of the race, and the use of the neural oracle is
key to doing so. Our results also suggest that a well-designed sketch can serve as a
regularizer. Due to constraints imposed by the sketch, the
policies for \torcs that \algo learns lead to smoother trajectories
than the corresponding neural policies, and can tolerate greater noise. The policies are also more easily transferred to new
domains, in particular race tracks not seen during training. Finally,
we show, using several properties, that the programmatic
policies that we discover are amenable to verification using
off-the-shelf symbolic techniques.

